%!TEX program = pdflatex
\documentclass[aps,%
twocolumn,%
floatfix,%
pra,%
amsfonts,%
longbibliography,%
groupedaddress,%
superscriptaddress]{revtex4-2}%

%%%%%%%%%%%%%%%%%%%%%%%%%%%%%%%%%%%%%%%%%%%%%%%%%%%%%%%%%%%%%%%%%%%%%%%%%
% Packages
%%%%%%%%%%%%%%%%%%%%%%%%%%%%%%%%%%%%%%%%%%%%%%%%%%%%%%%%%%%%%%%%%%%%%%%%%
\usepackage{amsmath}
\usepackage{amssymb}
\usepackage{amsthm}
\usepackage{bm,bbm} % to use the command \mathbbm{1}, \bm{x}
\usepackage{graphicx}
\usepackage{indentfirst}
\usepackage{enumitem}
\usepackage{tikz}
\usepackage{quantikz}
\usetikzlibrary{graphs} 
\usepackage[titletoc,toc,page,title]{appendix}
\definecolor{beamer@blendedblue}{rgb}{0.2,0.2,0.7}      % color used in beamer
\usepackage[bookmarks=false]{hyperref}
\hypersetup{colorlinks=true,%
citecolor=beamer@blendedblue,%
linkcolor=beamer@blendedblue,%
urlcolor=beamer@blendedblue,%
pdfstartview=FitH,%
bookmarksopen=true}

\def\UrlBreaks{\do\/\do-} % url break
\usepackage[T1]{fontenc}
\usepackage{times}
\usepackage{booktabs} % to use the \toprule \midrule command
%%%%%%%%%%%%%%%%%%%%%%%%%%%%%%%%%%%%%%%%%%%%%%%%%%%%%%%%%%%%%%%%%%%%%%%%%
% General Commands
%%%%%%%%%%%%%%%%%%%%%%%%%%%%%%%%%%%%%%%%%%%%%%%%%%%%%%%%%%%%%%%%%%%%%%%%%
\newtheorem{theorem}{Theorem}
\newtheorem{proposition}[theorem]{Proposition}
\newtheorem{lemma}{Lemma}
\newtheorem{corollary}{Corollary}
\newtheorem{algorithm}{Algorithm}
%%%%%%%%%%%%%%%%%%%%%%%%%%%%%%%%%%%%%%%%%%%%%%%%%%%%%%%%%%%%%%%%%%%%%%%%%
\usepackage{pifont}

%%%%%%%%%%%%%%%%%%%%%%%%%%%%%%%%%%%%%%%%%%%%%%%%%%%%%%%%%%%%%%%%%%%%%%%%%
% Quantum Commands
%%%%%%%%%%%%%%%%%%%%%%%%%%%%%%%%%%%%%%%%%%%%%%%%%%%%%%%%%%%%%%%%%%%%%%%%%
\DeclareMathOperator{\Tr}{Tr}
\DeclareMathOperator{\tr}{Tr}
\DeclareMathOperator{\NOT}{NOT}
\DeclareMathOperator{\diag}{diag}
\DeclareMathOperator{\rank}{rank}
\DeclareMathOperator{\conv}{conv}
\DeclareMathOperator{\sgn}{sgn}
\DeclareMathOperator{\Sp}{Span}
\DeclareMathOperator{\id}{id}
% \newcommand{\bra}[1]{\langle #1\rvert}
% \newcommand{\ket}[1]{\lvert #1\rangle}
% \newcommand{\proj}[1]{\lvert #1\rangle\!\langle #1\rvert}
\newcommand{\ox}{\otimes}
\newcommand{\1}{\mathbbm{1}}
\newcommand{\mean}[1]{\langle #1\rangle}
\newcommand{\gen}[1]{\langle #1\rangle}
\newcommand{\ketbra}[2]{\lvert #1\rangle\!\langle #2\rvert}
% \newcommand{\braket}[2]{\langle #1\vert #2\rangle}
%%%%%%%%%%%%%%%%%%%%%%%%%%%%%%%%%%%%%%%%%%%%%%%%%%%%%%%%%%%%%%%%%%%%%%%%%
%%% Mathematical symbols
%%%%%%%%%%%%%%%%%%%%%%%%%%%%%%%%%%%%%%%%%%%%%%%%%%%%%%%%%%%%%%%%%%%%%%%%%
\newcommand*{\cF}{\mathcal{F}}
\newcommand*{\cH}{\mathcal{H}}
\newcommand*{\cO}{\mathcal{O}}
\newcommand*{\cP}{\mathcal{P}}
\newcommand*{\cS}{\mathcal{S}}
\newcommand*{\bF}{\mathbb{F}}
\newcommand*{\bP}{\mathbb{P}}

% Colors
\definecolor{alizarin}{rgb}{0.82, 0.1, 0.26}
\definecolor{googleblue}{HTML}{4285F4}
\definecolor{googlered}{HTML}{DB4437}
\definecolor{googleyellow}{HTML}{F4B400}
\definecolor{googlegreen}{HTML}{0F9D58}
\definecolor{klevinblue}{HTML}{002FA7}
\definecolor{tiffanyblue}{HTML}{0ABAB5}

\newcommand{\SYC}[1]{\textcolor{blue}{(SYC: #1)}}

\begin{document}

\title{Entanglement induced complexity transition in high magic circuits}
%%%
\author{Siyuan Chen}
\affiliation{Hefei National Research Center for Physical Sciences at the Microscale and School of Physical Sciences, 
University of Science and Technology of China, Hefei 230026, China}%
%%%


\begin{abstract}
\noindent 
We propose a method to reduce arbitrary quantum circuit into non-adaptive measurement of certain graph state. 
% Result 1
For sampling amplitudes, the criterion involves determining whether percolation occurs on graphs. 
For Pauli expectation value estimation, the criterion evaluates whether a branching process can generate exponentially many descendants.
% Result 2
Both of the criteria can be simulated in polynomial time and related to operator stabilizer entropy of gadget operators.
% Conclusion
Our findings connect traditional simulation methods for Clifford + T circuits to Pauli path truncation method for noisy circuit and gives a general criterion in testing the supremacy in fault-tolerant quantum computation and NISQ quantum algorithm.
\end{abstract}
\date{\today}
\maketitle

%%%%%%%%%%%%%%%%%%%%%%%%%%%%%%%%%%%%%%%%%%%%%%%%%%%%%%%%%%%%%%%%%%%%%%%%%
%%%%%%%%%%%%%%%%%%%%%%%%%%%%%%%%%%%%%%%%%%%%%%%%%%%%%%%%%%%%%%%%%%%%%%%%%
%%%%%%%%%%%%%%%%%%%%%%%%%%%%%%%%%%%%%%%%%%%%%%%%%%%%%%%%%%%%%%%%%%%%%%%%%
\section{Introduction}
Fault-tolerant quantum computing aims to speed up tasks like quantum system simulation, large number factoring, optimization, and solving linear systems. 
Although fault-tolerant quantum computers are not yet realized, advances in quantum devices have enabled them to surpass classical supercomputers in certain tasks, particularly in the NISQ (noisy intermediate-scale quantum) era. 
In previous research, NISQ is defined as problems that can be solved by a polynomial-time probabilistic calssical algorithm with noisy quantum circuit that consist of (i). Preparation of $\text{poly}(n)$ noisy $\ket{0}$ qubits, (ii). Application of $\text{poly}(n)$ layers' noisy two-qubit gates, and (iii).Noisy computational basis measurements~\cite{chen2023complexity}. 
A concrete example is the depth-$d$ circuit with local noise after each layer. This circuit output quantum state
\begin{align}
    C(\rho_0) = \mathcal{N}_0\{U_1\cdots \mathcal{N}_{d-1}\{ U_d\mathcal{N}_d\{\rho_0\}U_d^\dagger\}\cdots U_1^\dagger \},
\end{align}
where $\mathcal{D}$ is the composition of local channel. Then measurement on the computational basis is performed, useful information such as $\langle Z^{s} \rangle$ or $\bra{x}C(\rho_0)\ket{x}$ is extracted.
 \begin{figure}[htbp!]
     \centering
     \includegraphics[width=1\linewidth]{pic/NISQ_circuit.png}
     \caption{Demonstration of NISQ circuit. }
     \label{fig:NISQ-circuit}
 \end{figure}

A line of research has focused on the computational power of NISQ circuits. This understanding could help identify the key factors necessary for achieving quantum advantage. Thus many simulation algorithm were proposed for NISQ circuits. An upper bound have been proved by Aharonov \emph{et.al.} ~\cite{aharonov1996limitations}, which states that NISQ circuit has the same computational power of logarithmic depth quantum circuit. This bound suggests us to focus on algorithms for NISQ circuits with fixed structure. 

The first fraction of structured NISQ circuit is ensembles of NISQ circuit with certain property such as ensemble with hiding property and anticoncentration~\cite{aharonov2023polynomial} , low-average ensembles~\cite{schuster2024polynomial} and variational quantum circuit with uniform random parameters~\cite{shao2023simulating}. These properties of ensemble are naturally satisfied by many random circuit sampling experiments. For those circuits, Pauli-path concatenation method can be applied to implement the \textit{mean value estimation task} and \textit{sampling task}, both within constant additive error and polynomial time. 
However, the introduction of circuit ensemble might overwhelmed the effect of noise because it is recently show that for noiseless scrambling circuit ensemble, Pauli-path concatenation method can also run in polynomial time and estimates mean values within additive error on average\cite{angrisani2024classically}. A more subtle problem of this method is the ineffectiveness in the case of non-unital channel. Non-unital noise such as amplitude damping channel takes the main part of the noise in neutral atoms and ion-trap quantum computers~\cite{cong2022hardware,sahay2023high} . This calls for new simulation methods in order to benchmark Random Circuit Sampling (RCS) tasks and Variational Quantum Algorithms (VQAs).

Another fraction of NISQ circuit are fixed quantum circuit with bounded quantum resource. 
Resource theories characterize the classically tractable operations in a certain framework. In this work, we particularly focus on entanglement and magic.
There are many well-known algorithms for near-stabilizer quantum circuits ~\cite{bravyi2016improved,bravyi2019simulation,seddon2021quantifying} and 
low entanglement quantum circuits~\cite{markov2008simulating,van2007classical}. 
The relation between magic and entanglement also arose many attention and recently~\cite{gu2024magic,masot2024stabilizer}.
It turns out that integrate these two approaches will present efficient simulation method. However, it is still unclear how the existence of noise will break the power of quantum resource in fixed quantum circuits. 
It is expect to show a trade off between the noise and quantum resource. Guillermo \emph{et.al.} move the fist step and show that  if noise is introduced faster than T gates, the Pauli-path concatenation method still works ~\cite{gonzalez2024pauli}. 
Recently, efficient algorithms beyond Pauli-path concatenation method are also proposed when noise exist on fixed Clifford-fixed circuits~\cite{nelson2024polynomial} and fixed IQP circuits~\cite{rajakumar2024polynomialtimeclassicalsimulationnoisy}. They both based on percolation of certain graphs and it is still unclear how this effect related to the magic or entanglement resource in the circuit.

In this work, we proposed a unified simulation method for fixed NISQ circuit with bounded quantum resource, with respect to the task of sampling amplitudes.
The algorithms for sampling amplitudes connect T-gadget simulation method to the percolation based algorithm, whose efficiency are bounded by entanglement measure in terms of Schmidt-rank width~\cite{van2007classical}. The threshold of noise is related to the phase-transition tree width of random graph states.
Noted that our algorithms still work in the case of non-unital noise, as long as magic is not introduced by noise channel~\cite{turkeshi2024coherent}, which could be used as a standard criterion for all quantum circuit on NISQ-era. 


\section{Background}
\subsection{Percolation method}
Recently, simulation method for noisy quantum circuit percolation 

\subsection{Gadget-based simulation}

Consider a quantum circuit consisting of \( t \) \( T \)-gates \((T = \text{diag}(1, e^{i\pi/4}))\) and \( c \) Clifford gates. It is well known that such a circuit can be simulated using a gadget-based simulation method. In this approach, the \( T \)-gates are replaced with multiple \( T \)-gate gadgets. The gadget circuit consists of (i). $t$ ancillary qubits in the product state \(\ket{H} = H\ket{T} = \cos(\pi/8)\ket{0} + \sin(\pi/8)\ket{1}\),(ii). parallel \( CZ \)-gates that inject magic into the original circuit, and (iii).measurement of the ancillary qubits in the \(\{\ket{+}, \ket{-}\}\) basis.

\begin{figure}[!htbp]
    \centering
    \includegraphics[width=1\linewidth]{pic/gadget_based.png}
    \caption{Illustration of our gadget-based noisy circuit simulation. T-gates at each depth is firstly replaced by post-selected T-gadget. }
    \label{fig:gadget_based}
\end{figure}
If the measurement outcome for an ancillary qubit is \(\ket{-}\), an \( S \)-gate is applied to the corresponding input qubit.

The gadget circuit replicates the same distribution on the input qubits as the original circuit. By replacing the measurement of the ancillary qubits with projectors onto \(\ket{+}\), the \( S \)-gate corrections are eliminated, yielding the following expression
\begin{align}
    p(x) = 2^t \left|\bra{+^{\otimes n} \otimes A^{\otimes t}} V^\dagger \ket{x_+^{\otimes n} \otimes +^{\otimes t}}\right|^2.
\end{align}
Starting from the stabilizer of the state \(\ket{x_+^{\otimes n} \otimes A^{\otimes t}}\), and assuming the action of a single Clifford gate on Pauli operators can be computed in \( O(1) \) time, the result of each generator acted upon by \( V \) can be computed in \( O(c + t) \) time, where \( c \) is the number of Clifford gates. Furthermore, the \( w \)-qubit marginal distribution of interest, involving \( w \) qubits, can be computed in \( O((w + t)(c + t)) \) time. The corresponding subspace of stabilizer generators is denoted as \(\mathcal{V}_0\), and the probability expression simplifies to:
\[
p(x) = 2^t \bra{+^{\otimes n}  A^{\otimes t}} \frac{\sum_{g \in \mathcal{V}_0} g}{2^{w + t}} \ket{+^{\otimes n}  A^{\otimes t}}.
\]

If any \( Z \) or \( Y \) operators appear on the first \( n \) qubits, the result is \( 0 \). Therefore, we extract from \(\mathcal{V}_0\) only those elements that act as \( I \) or \( X \) on the first \( n \) qubits. These elements form a commuting subspace \(\mathcal{V}_1\). Using Gaussian elimination, the generators of \(\mathcal{V}_1\) can be determined. If \(-I\) is among the generators, then \( p(x_w) = 0 \); otherwise, we find the stabilizer subspace associated with these generators and compute:
\[
p(x) = 2^t \bra{A^{\otimes t}} \frac{\sum_{g \in \mathcal{V}_1} g}{2^{w + t}} \ket{A^{\otimes t}}.
\]

The second part of the algorithm involves decomposing the magic state \(\ket{A^{\otimes t}}\), identifying \(\ket{\psi_1}, \ldots, \ket{\psi_{\chi(t)}} \in \mathcal{S}\) such that:
\[
\ket{A^{\otimes t}} = \sum_{a=1}^{\chi(t)} c_a \ket{\psi_a}.
\]
Finding such a decomposition, including the coefficients \(\{c_a, \ket{\psi_a}\}\), can be achieved in polynomial time. The stabilizer rank \(\chi(t)\) of \(\ket{A^{\otimes t}}\) is conjectured to grow exponentially with \( t \), although a rigorous lower bound has yet to be proven~\cite{mehraban2024quadratic}. Once the decomposition is complete, one needs to compute the overlaps of \(\chi(t)^2\) stabilizer states:
\[
p_a(x) = 2^t c_a c_b \bra{\psi_a} \frac{\sum_{g \in \mathcal{V}_1} g}{2^{w + t}} \ket{\psi_b},
\]
and sum these to obtain \( p(x) \). 

Using sampling and approximate stabilizer rank decomposition, this approach can be extended to be more efficient~\cite{bravyi2016improved}. A 
%%%%%%%%%%%%%%%%%%%%%%%%%%%%%%%%%%%%%%%%%%%%%%%%%%%%%%%%%%%%%%%%%%%%%%%%%


%%%%%%%%%%%%%%%%%%%%%%%%%%%%%%%%%%%%%%%%%%%%%%%%%%%%%%%%%%%%%%%%%%%%%%%%%
\section{Summary of main results}
In traditional gadget-based methods, although completely stabilizer-preserving(CSP) channels are considered, it still leads to exponential growth of simulation complexity depending on the number of non-Clifford gates. 
This means that they can only applied to near-Clifford circuit with $\ln n$ numbers of non-Clifford gates.
However, it has been proven that noisy circuits with polynomial numbers but sparse T-gates can be efficiently simulated by Pauli path truncation method. 
On the other hand, the effectiveness of Pauli path truncation algorithm stem from the property of Pauli unital noise, which exponentially suppress high-weight Pauli paths while not spreading them. 
This prevent the application of truncation algorithm on more general noise.
Nevertheless, many non-unital qubit channels, though exponentially spread Pauli operators, will not introduce magic resource at all. Near-Clifford circuit with such noise can be efficiently simulated~\cite{seddon2021quantifying,}
Therefore, a natural question arises:
\textit{Can we generalized the gadget simulation method such that it can be applied to 'sparse' magic circuit with  arbitrary CSP channel?}

Motivated by these questions, we propose three randomized algorithms for the simulation of noisy circuits based on the gadget method, for both Pauli averages approximation task and amplitude sampling tasks. This is concludes in

Given a fixed noisy Clifford+T circuit, our method provides a criterion to test the 'sparseness' of T gates in the circuit, such that the impact of noise obscures the simulation complexity.  
Once the criterion yields an affirmative conclusion, we can invoke the corresponding algorithms to estimate the circuit's sampling amplitude and Pauli expectation values.  


\begin{theorem}
    If the randomized graph generated by Algorithm 1 can be split into $V_1,\cdots,V_m$ that satisfied
    \begin{align}
        \text{Prob}(\max_i|V_i| > x) \leq e^{-\alpha x}
    \end{align}
    Then there exist a Monte-Carlo method that samples from $Q_C$ such that $\|Q_C - P_C\|_{\text{TVD}} \leq \varepsilon$ within time $T \leq O(n^{c \alpha^{-1}+c'})$ . 
\end{theorem}


%%%%%%%%%%%%%%%%%%%%%%%%%%%%%%%%%%%%%%%%%%%%%%%%%%%%%%%%%%%%%%%%%%%%%%%%%
%%%%%%%%%%%%%%%%%%%%%%%%%%%%%%%%%%%%%%%%%%%%%%%%%%%%%%%%%%%%%%%%%%%%%%%%%
\section{ Single qubit incoherent noise}
Any single-qubit noise channel can be expressed as
\[
\mathcal{N}(\rho) = U \Lambda(V \rho V^\dagger)U^\dagger,
\]
where $U,V$ are unitary gates that represent the coherent component of noise channels and \(\Lambda\left(\frac{1}{2}[\mathbb{I} + \mathbf{w} \cdot \vec \sigma]\right) = \frac{1}{2}[\mathbb{I} + (\mathbf{t} + \Gamma\mathbf{w}) \cdot  \vec\sigma]\) is the incoherent parts, with \(\Gamma\) a diagonal matrix with diagonal entries less than 1.
Two component has different impact on the complexity of circuit. It is recently aware that coherent errors can generate high-magic resource states for universal quantum computing. ~\cite{niroula2024phase,turkeshi2024coherent}
To prevent the coherent parts introducing magic to the circuit, we only discuss the case where $U,V$ are both Clifford gates. The effect of Clifford coherent noise can be further included into the Clifford layers. Thus it is sufficient to consider noise channel in form of $\Lambda$. 
For amplitude damping channel $\Lambda(\rho) = K_0 \rho K_0^\dagger + K_1 \rho K_1^\dagger$, where $K_0 = \text{diag}(1, \sqrt{1 - p}), K_1 = [0,\sqrt{q}; 0,0]$ . 

%%%%%%%%%%%%%%%%%%%%%%%%%%%%%%%%%%%%%%%%%%%%%%%%%%%%%%%%%%%%%%%%%%%%%%%%%
%%%%%%%%%%%%%%%%%%%%%%%%%%%%%%%%%%%%%%%%%%%%%%%%%%%%%%%%%%%%%%%%%%%%%%%%%
%%%%%%%%%%%%%%%%%%%%%%%%%%%%%%%%%%%%%%%%%%%%%%%%%%%%%%%%%%%%%%%%%%%%%%%%%
\section{Algorithm for amplitude sampling}
\begin{theorem}[Noisy reduction to graph states]
\label{theo:Noisy_reduction}
 $\mathcal{N}$ is incoherent single-qubit channel that has the form
 \begin{align}
    \mathcal{N}\left(\frac{1}{2}[\mathbb{I} + \mathbf{w} \cdot \vec \sigma]\right) = \frac{1}{2}[\mathbb{I} + (\mathbf{t} + \Gamma\mathbf{w}) \cdot  \vec\sigma] .
 \end{align}
In the absence of noise, 
Clifford operation $\text{Cl}$ on stabilizer state $\ket{S}$ can be decomposed into a set of control-Z gates $G$  followed by local Clifford gates $L$  as if the noise do not exist.
\begin{align}
    &\text{Cl} \ket{S}= L \cdot G \ket{S}, \\
    &\text{Cl} \circ \mathcal{N}[\ket{S}\!\bra{S}] = LG \circ \mathcal{N}[\ket{S}\!\bra{S}]
\end{align}
\end{theorem}


%%%%%%%%%%%%%%%%%%%%%%%%%%%%%%%%%%%%%%%%%%%%%%%%%%%%%%%%%%%%%%%%%%%%%%%%%
%%%%%%%%%%%%%%%%%%%%%%%%%%%%%%%%%%%%%%%%%%%%%%%%%%%%%%%%%%%%%%%%%%%%%%%%%
\subsection{subsection1}



%%%%%%%%%%%%%%%%%%%%%%%%%%%%%%%%%%%%%%%%%%%%%%%%%%%%%%%%%%%%%%%%%%%%%%%%%
%%%%%%%%%%%%%%%%%%%%%%%%%%%%%%%%%%%%%%%%%%%%%%%%%%%%%%%%%%%%%%%%%%%%%%%%%
%%%%%%%%%%%%%%%%%%%%%%%%%%%%%%%%%%%%%%%%%%%%%%%%%%%%%%%%%%%%%%%%%%%%%%%%%
\section{section2}


%%%%%%%%%%%%%%%%%%%%%%%%%%%%%%%%%%%%%%%%%%%%%%%%%%%%%%%%%%%%%%%%%%%%%%%%%
%%%%%%%%%%%%%%%%%%%%%%%%%%%%%%%%%%%%%%%%%%%%%%%%%%%%%%%%%%%%%%%%%%%%%%%%%
\subsection{subsection}



%%%%%%%%%%%%%%%%%%%%%%%%%%%%%%%%%%%%%%%%%%%%%%%%%%%%%%%%%%%%%%%%%%%%%%%%%
%%%%%%%%%%%%%%%%%%%%%%%%%%%%%%%%%%%%%%%%%%%%%%%%%%%%%%%%%%%%%%%%%%%%%%%%%
%%%%%%%%%%%%%%%%%%%%%%%%%%%%%%%%%%%%%%%%%%%%%%%%%%%%%%%%%%%%%%%%%%%%%%%%%
\section{Conclusions}



%%%%%%%%%%%%%%%%%%%%%%%%%%%%%%%%%%%%%%%%%%%%%%%%%%%%%%%%%%%%%%%%%%%%%%%%%
% Ack
%%%%%%%%%%%%%%%%%%%%%%%%%%%%%%%%%%%%%%%%%%%%%%%%%%%%%%%%%%%%%%%%%%%%%%%%%
\section*{Acknowledgements} 



%%%%%%%%%%%%%%%%%%%%%%%%%%%%%%%%%%%%%%%%%%%%%%%%%%%%%%%%%%%%%%%%%%%%%%%%%
% References
%%%%%%%%%%%%%%%%%%%%%%%%%%%%%%%%%%%%%%%%%%%%%%%%%%%%%%%%%%%%%%%%%%%%%%%%%

\bibliographystyle{apsrev4-2}
\bibliography{references.bib}
%apsrev4-2.bst 2019-01-14 (MD) hand-edited version of apsrev4-1.bst
%Control: key (0)
%Control: author (8) initials jnrlst
%Control: editor formatted (1) identically to author
%Control: production of article title (0) allowed
%Control: page (0) single
%Control: year (1) truncated
%Control: production of eprint (0) enabled

%%%%%%%%%%%%%%%%%%%%%%%%%%%%%%%%%%%%%%%%%%%%%%%%%%%%%%%%%%%%%%%%%%%%%%%%%%%%%%%%%%%%%%%%%%%%%%%%%%%
% Appendices
%%%%%%%%%%%%%%%%%%%%%%%%%%%%%%%%%%%%%%%%%%%%%%%%%%%%%%%%%%%%%%%%%%%%%%%%%%%%%%%%%%%%%%%%%%%%%%%%%%%
% \appendix
% \onecolumngrid
% https://tex.stackexchange.com/questions/488302/referencing-an-appendix-while-using-prl-style
\makeatletter
\newcommand{\appendixtitle}[1]{\gdef\@title{#1}}
\newcommand{\appendixauthor}[1]{\gdef\@author{#1}}
\newcommand{\appendixaffiliation}[1]{\gdef\@affiliation{#1}}
\newcommand{\appendixdate}[1]{\gdef\@date{#1}}
\makeatother

% Define the custom title block
\makeatletter%
\newcommand{\appendixmaketitle}{%
\begin{center}%
\vspace{0.4in}%
{\Large \@title \par}%
% {\normalsize
%   \lineskip .5em%
%   \begin{tabular}[t]{c}%
%     \@author
%   \end{tabular}\par}%
% \vspace{5pt}
% {\itshape \@affiliation \par}%
\end{center}%
\par%
}%
\makeatother%

\setcounter{secnumdepth}{2}
\appendix
\widetext
\newpage

\appendixtitle{\bf 
Supplemental Material for "''}
% \appendixauthor{Siyuan Chen, Wei Xie, and Kun Wang} 
% \appendixaffiliation{
% Institute for Quantum Computing, Baidu Research, Beijing 100193, China \\
% Physics Department, Jilin University \\
% School of Computer Science and Technology, University of Science and Technology of China}
\appendixmaketitle
\vspace{0.2in}

%%%%%%%%%%%%%%%%%%%%%%%%%%%%%%%%%%%%%%%%%%%%%%%%%%%%%%%%%%%%%%%%%%%%%%%%%%%%%%%%%%%%%%%%%%%%%%%%%%%
%%%%%%%%%%%%%%%%%%%%%%%%%%%%%%%%%%%%%%%%%%%%%%%%%%%%%%%%%%%%%%%%%%%%%%%%%%%%%%%%%%%%%%%%%%%%%%%%%%%
%%%%%%%%%%%%%%%%%%%%%%%%%%%%%%%%%%%%%%%%%%%%%%%%%%%%%%%%%%%%%%%%%%%%%%%%%%%%%%%%%%%%%%%%%%%%%%%%%%%
The contents of the supplementary material are structured as follows: 
In Appendix..
%%%%%%%%%%%%%%%%%%%%%%%%%%%%%%%%%%%%%%%%%%%%%%%%%%%%%%%%%%%%%%%%%%%%%%%%%%%%%%%%%%%%%%%%%%%%%%%%%%%
%%%%%%%%%%%%%%%%%%%%%%%%%%%%%%%%%%%%%%%%%%%%%%%%%%%%%%%%%%%%%%%%%%%%%%%%%%%%%%%%%%%%%%%%%%%%%%%%%%%
%%%%%%%%%%%%%%%%%%%%%%%%%%%%%%%%%%%%%%%%%%%%%%%%%%%%%%%%%%%%%%%%%%%%%%%%%%%%%%%%%%%%%%%%%%%%%%%%%%%

\section{Noise channels}
In this appendix, we show some critial properties of single qubit quantum channels. It will been shown that the incoherent parts of single qubit noise channels has a great inpact on circuit complexity.
\subsection{Diagonalization property}


\subsection{Graph-breaking property}
It is possible to back-propagate the stabilizer state before the incoherent channel. This is the stabilizer state that generated through measuring Pauli operators $\{\sigma_1,\sigma_2,\sigma_3\}$. 

\section{Graphical compilation of noisy quantum circuit}
It appears that any polynomial-depth quantum circuit can be compiled into a sequence of non-adaptive T-state measurements on a graph state. This compilation simplifies quantum error correction and quantum computation. For instance, existing hardware would only need to repeatedly perform single-qubit T measurements on a graph state to obtain randomness, which can then be processed classically to synthesize the randomness required for any polynomial-depth circuit.

\subsection{Graphical compilation for amplitude estimation}
Consider a depth-$d$ noisy quantum circuit of $n$ qubits, as manifest in Figure~\ref{fig:NISQ-circuit}. Here $\mathcal{N}_i$ is Pauli noise and $\mathcal{C}_i, \mathcal{T}_i$ are Clifford layers and T-gates layers.  In the classical simulation algorithm, we first consider the approximation of fixed sampling amplitude
\begin{align}
    p_C(+^{\otimes n}) = \text{tr}\left(\ket{+}\!\bra{+}^{\otimes n} \cdot\mathcal{N}_0  \circ (\mathcal{C}_1\circ \mathcal{T}_1 \circ\mathcal{N}_1)\circ \cdots\circ(\mathcal{C}_d\circ \mathcal{T}_d \circ \mathcal{N}_d)\left[\ket{0}\!\bra{0}^{\otimes n} \right]\right).
\end{align}

According to the T-gadget simulation method, as mentioned in last section, all T-gates layers can be replaced by magic state inputs $\ket{A} =e^{-i\frac{\pi}{8}}H\ket{T} = \cos {\pi}/{8}\ket{0} + \sin{ \pi /8} \ket{1}$ , which then left circuit with Clifford gates and single qubit noise channels only. The amplitude of post-selected T-gadget circuit relates to the original circuit sampling amplitude through
\begin{align}
    p_C(+^{\otimes n}) = 2^{t} \cdot \text{tr}\left(\ket{+}\!\bra{+}^{\otimes (n+t)} \cdot\mathcal{N}_0  \circ (\mathcal{C}_1\circ \mathcal{CZ}_1 \circ\mathcal{N}_1)\circ \cdots\circ(\mathcal{C}_d\circ \mathcal{CZ}_d \circ \mathcal{N}_d)\left[\ket{0}\!\bra{0}^{\otimes n}  \otimes \ket{A}\!\bra{A}^{\otimes t} \right]\right),
\end{align}
where $t$ is the total number of $T$ gates and $\mathcal{CZ}_i$ is control-Z gates layers that was introduced due to gadget replacement. 
In this circuit, the magic state evlove through Clifford gates and then measured on stabilizer state $\rho=\ket{+}\!\bra{+}^{\otimes n}$, along with post-selection states $\ket{+}^{\otimes t}$.
It is valueable to regards state evolution under Heisenberg point of view. The measurement basis can be moved foreward through Clifford gates and incoherent noise, and finally generates operator $\hat O$ .
\begin{align}
\label{eq:conjugated_amplitude}
    p_C(+^{\otimes n})   &= 2^{t} \cdot \text{tr}\left(\ket{0A}\!\bra{0A}^{\otimes (n+t)}  \cdot (\mathcal{N}_1^\dagger\circ \mathcal{CZ}_1^\dagger\circ \mathcal{C}_1^\dagger)\circ \cdots\circ(\mathcal{N}_d^\dagger\circ \mathcal{CZ}_d^\dagger\circ \mathcal{C}_d^\dagger)\circ \mathcal{N}_0^\dagger\left[ \ket{+}\!\bra{+}^{\otimes (n+t)} \right]\right)\\
    &:=2^{t} \cdot \text{tr}\left(\ket{0A}\!\bra{0A}^{\otimes (n+t)}  \cdot \hat O\right).
\end{align}
We then define $\hat O' = 2^{n}\text{tr}_{\{1,\cdots n\}} (\hat O \ket{0}\!\bra{0}^{\otimes n})$ and $p_C(+^{\otimes n}) = 2^{(t-n)}\text{tr}(\hat O ' \ket{A}\!\bra{A}^{\otimes t})$. In the following, we will show that the operator $\hat O'$ can be simulated through sampling different graph states. The property of this graph ensemble reflects the hardness of this quantum circuit.
%%%%%%%%%%%%%%%%%%%%%%%%%
%%%%%%%%%%%%%%%%%%%%%%%%%%%%%
%%%%%%%%%%%%%%%%%%%%%%%%%%%%%%%%%%

\subsubsection{Noiseless case}
For noiseless Clifford+ T circuits, $\mathcal{CZ}^\dagger,\mathcal{C}^\dagger$ is a valid Clifford operation. Thus the operator $\hat O'$ is a valid stabilizer state $\ket{S}\!\bra{S}$. Thanks to the Gottesman-Knill theorem, the operator $\hat O'$ can be efficiently simulated through tracing the stabilizer state $\ket{S}$ over $n +t$ qubits. Then the propability can be calulated through
\begin{align}
    p_c(+^{\otimes n}) &={2^{t}} \left|\langle S | 0A \rangle^{\otimes (n+t)}\right|^2\\
    &={2^{t-n}} \left|\langle S' | A \rangle^{\otimes t}\right|^2 ={2^{t-n}} \left|\langle G|L^\dagger| A \rangle^{\otimes t}\right|^2
\end{align}
where $\ket{S'}$ is a stabilizer state on $t$ qubits that generated through measuring the first $n$ qubits of $\ket{S}$ state into $\ket{0}$ basis. To further transform $\ket{S'}$ into graph state $\ket{G}$ in the last line, we refer the lemma below~\cite{van2004graphical}.
\begin{lemma}\label{lemma: stabilizer_to_graph}
   Given a $n$-qubit stabilizer state $\ket{S'}$ , a local Clifford operation $L$ can be find in $\text{poly}(n)$ time such that stabilizer state $\ket{S'} = L \ket{G}$, where $\ket{G}$ is a graph state with graph $G = (V,E)$.  
\end{lemma}

Now we compile the amplitude  $p_C(+^n)$ into a $t$ vertex graph and corresponding single qubit measurement on basis $L_i^\dagger\ket{A}$.  
For the same circuit, the other amplitudes can be compiled in the same procedures and will give a same graph $G$ but a different single qubit measurements $L_i^{'\dagger} \ket{A}$. 
For arbitrary graph state, it is believed hard to classically approximate its overlap between single qubit magic state within $O(2^{-t})$ error. 
We noted that it is possible to estimate this value within an constant error, because $\ket{S}$ and $\ket{0A}^{\otimes (n+t)}$ are both tractable states~\cite{nest2009simulating}. 
However, for certain families of stabilizer states, an efficient classical description exist. 

Previous research on measurement-based quantum computation(MQC) shows that, for the family of graph states with $r =O(\text{log~} n)$ graph rank width, we can first efficiently find its tree tensor network(TNN) representation with bounded tensor dimension $D =O(2^r) = \text{poly}(n)$ and then calculate this overlap accurately~\cite{van2007classical}. Thus it is valuable to test the rank-width of graph $G$ to ensure the hardness of estimating probablity amplitude.

%%%%%%%%%%%%%%%%%%%%%%%%%%%%
%%%%%%%%%%%%%%%%%%%%%%%%%%%%%%
%%%%%%%%%%%%%%%%%%%%%%%%%%%%%%%%%%
\subsubsection{Diagonalization in noisy case}
Now we consider quantum circuit with single qubit incoherent noise. Similar to Gottesman-Knill theorem in the noiseless case, it is also possible to 'back-propagates' stabilizer state through operations and reduce the stabilizer state into graph states. 
We first show that lemma~\ref{lemma: stabilizer_to_graph} is still correct as if the noise does not exist. This is showed by theorem~\ref{theo:Noisy_reduction} below.
\begin{theorem}[Noisy reduction to graph states]
\label{theo:Noisy_reduction}
Use $\text{Cl}$ to represent Clifford gate and $\mathcal{Cl}[\rho]:= \text{Cl}\rho \text{Cl}^\dagger$ to imply unitary operation on mixed state.
 $\mathcal{N}$ is Pauli noise channel.
 
 % that has the form
 % \begin{align}
 %    \mathcal{N}\left(\frac{1}{2}[\mathbb{I} + \mathbf{w} \cdot \vec \sigma]\right) = \frac{1}{2}[\mathbb{I} + (\mathbf{t} + \Gamma\mathbf{w}) \cdot  \vec\sigma] .
 % \end{align}
In the absence of noise, Clifford operation on stabilizer state $\ket{S}$ can be equivalently realized by a set of $CZ$ gates $G$  followed by local Clifford gates $L$ as $\text{Cl} \ket{S}= LG \ket{S}.$ This equation is still satisfied under existence of noise $\mathcal{N}$, which means $\mathcal{Cl} \circ \mathcal{N}[\ket{S}\!\bra{S}] = LG \circ \mathcal{N}[\ket{S}\!\bra{S}]$. 
\end{theorem}
\begin{proof}
    We first prove the existence of the decomposition in the noiseless case. First note the fact the every stabilizer state is local equivalence to graph states. Then $\ket{S} = L_SG_S \ket{+}$. We conclude that
    \begin{align}
        \text{Cl} \ket{S} &= L' G' \ket{+} = L' G'G_S L_S^\dagger L_SG_S \ket{+} \\
        &=L' G^{''} L_S^\dagger \ket{S} = L'L_S^\dagger L_S G^{''} L_S^\dagger \ket{S}=L^{'''} G^{'''} \ket{S}
    \end{align}
    The first equality use graph state decomposition of Clifford state $\text{Cl}\ket{S}$. The last equality is valid because $L_S G^{''} L_S^\dagger$ conjugates $CZ$ gate into control-Pauli gates, which can be immediately replaced by multiple $CZ$ gates and commutable single qubit phase gates, as illustrated by Figure~\ref{theo:Noisy_reduction} . Thus the operation of $L_S G^{''} L_S^\dagger$ are still multiple $CZ$ gates, which is denoted as $G^{'''}$ in the equality.
    Given the existence of $LG$ decomposition with respect to stabilizer state $\ket{S}$, Let $Cl = LAG$. Using the relation $LG \ket{S} = Cl \ket{S}$, it follows that $A G \ket{S} = G \ket{S}$. Thus, $G^\dagger A G$, as an element in Clifford group, is a stabilizer of the stabilizer state $\ket{S}$ and can be expressed as a product of Pauli operators. The following transformations then applied:
\begin{align}
\text{Cl} \circ \mathcal{N}[\ket{S}\bra{S}] &= \text{Cl} \mathcal{N} \text{Cl}^\dagger \circ \text{Cl}[\ket{S}\bra{S}] \\
&= \text{Cl} \mathcal{N} \text{Cl}^\dagger \circ LG[\ket{S}\bra{S}] \\
&= LG \circ (G^\dagger L^\dagger \text{Cl}) \mathcal{N} (\text{Cl}^\dagger LG)[\ket{S}\bra{S}] \\
&= LG \circ (G^\dagger AG) \mathcal{N} (G^\dagger A^\dagger G)[\ket{S}\bra{S}]\\
&= LG \circ  \mathcal{N}[\ket{S}\bra{S}].
\end{align}
In the last equality, we use the fact that $G^\dagger A G$ is only a product of Pauli operators, which has no effect on the incoherent noise channel, thus $(G^\dagger AG) \mathcal{N} (G^\dagger A^\dagger G) = \mathcal{N}$.
\end{proof}
For example, the control-Pauli gates on graph state with Pauli noise can be simplified. In the noiseless case, control-Pauli operations on graph state can be efficiently transformed into equivalent multiple control-Z gates according to the rule in Fig.~\ref{fig:CP_CZ_reduction}~\cite{mor2023noisy}. 
lemma~\ref{theo:Noisy_reduction} tell us these rules are still satisfied even if the graph states are suffering incoherent noise. 
\begin{figure}
    \centering
    \includegraphics[width=0.5\linewidth]{pic/CP_to_CZ.png}
    \caption{The control-X, control-Y and control-Z operation on arbitrary graph state $\ket{G}, G = (V,E)$.  Consider control-$P$ operation with control qubit $i$ and target qubit $j$. We use $CZ$ gate on $(i,i)$ means $\sigma_3$ gates on qubit corresponding to the  $i$-th vertex. When $P=\sigma_1$ ,it is equivalent to $CZ$ gates on pair $\{(i,k)|(k,j) \in E\}$. When $P=\sigma_2$ ,it is equivalent to $CZ$ gates on pair $\{(i,j)\}\cup\{(i,k)|(k,j) \in E\}$ , with a local Clifford correction $S$ on control qubit. }
    \label{fig:CP_CZ_reduction}
\end{figure}
A direct application of theorem~\ref{theo:Noisy_reduction} is to reduce the Clifford operation inside  $\mathcal{N}_i^\dagger\circ \mathcal{CZ}_i^\dagger\circ \mathcal{C}_i^\dagger$ in Eq.~\eqref{eq:conjugated_amplitude} into $CZ$ gates and local Clifford operations $L$ while not change the single qubit incoherent noise. Use this reduction, we can reduce the whole conjugated operation in Eq.~\eqref{eq:conjugated_amplitude} , except the noise layers, into a noisy IQP circuits, as illustrated in Figure~\ref{fig:Percolated_stabilizer_state}(d). During the process, we use the rule that mentioned in Fig.~\ref{fig:CP_CZ_reduction}~\cite{mor2023noisy} to transform control-Pauli operations on graph state into equivalent multiple control-Z gates. Local Clifford gates shift Pauli noise $\mathcal{N}_i^\dagger$ to $\mathcal{N'}_{i}^\dagger$, which is also a incoherent noise and can be efficiently computed. 
We call this process as\textit{ diagonalization of Pauli-noise Clifford circuit}. The relation between the diagonalized circuit with the original circuits are concluded and proved in theorem~\ref{theo:diagonalzation} below.
\begin{figure}[htbp!]
    \centering
    \includegraphics[width=1\linewidth]{pic/diagonalization.png}
    \caption{Illustration of our gadget-based noisy circuit simulation. T-gates at each depth is firstly replaced by post-selected T-gadget, as shown in Fig~\ref{fig:gadget_based}~(b). Then noise channels are moved forward to the begining. Noted this operation can be efficiently simulated as long as noise channel are mixed by $\text{Poly}(n)$ Clifford gates. In the end, as shown in Fig~\ref{fig:gadget_based}~(c), those noisy channel reduce the magic resource of the input magic state $\ket{0}^{\otimes n} \ket{T}^{\otimes t_1}\cdots\ket{T}^{\otimes t_d}$, which might induce an efficient gadget-based simulation strategy.  }
    \label{fig:Percolated_stabilizer_state}
\end{figure}



\begin{theorem}[Diagonalzation of Pauli-noise Clifford circuit]
\label{theo:diagonalzation}
Consider NISQ circuit with $d$ layer of Clifford gates  $\{\text{Cl}_i\}_{i =1}^d$  and $T$ gates. After replacing $T$ gates with $T$-gadget and decompose Clifford operation $\{\text{Cl}_i\}_{i =1}^d$ into $CZ$ gates $\{G_i\}_{i =1}^d$ as proved in Lemma~\ref{lemma:Noisy_reduction} , we get the circuit in Figure~\ref{fig:Percolated_stabilizer_state}(b). Those transformed gates and noisy channel can be calculated from NISQ circuit through the relation below
\begin{align}
    &G_j G_{j-1}\cdots G_1=D_j^\dagger,
    \text{where~}\text{Cl}_j^\dagger \cdots \text{Cl}_1^\dagger \ket{+}^{\otimes n} = L_j^\dagger D_j^\dagger \ket{+}^{\otimes n}   ~\forall j = 1, \cdots, d.  \\
     &P_{j}= A_j Z A_j^\dagger,\text{where~} A_d = L_d, A_j =   (G_{j+1} \cdots G_d )L_d(\text{Cl}_{j+1} \cdots\text{Cl}_d)^\dagger ~\forall j = 1, \cdots,d-1. \\
    &\mathcal{N}_j^{'}(\rho) = L_j\circ \mathcal{N}_j \circ L_j^\dagger(\rho) .
\end{align}
Here the $P_j$ is the target operator of control-Pauli gates and should be further transformed into multiple control-Z gates according to the rule in Fig.~\ref{fig:CP_CZ_reduction} . 
\end{theorem}
\begin{proof}
   The form of $P_i$ can be easily verified through replacing $P_j = A_j Z A_j^\dagger$ in the diagonalized circuit and use the fact that $CP_u = A_j (CZ) A_j^\dagger$. This will give the original quantum circuit, except that the Pauli noise channel $\mathcal{N}_{j-1}^{'}$ after $G_j$ are conjugated as $A_{j-1}^\dagger\circ \mathcal{N}_{j-1}^{'}\circ A_{j-1} $. With input state $\rho$, this channel operates as
   \begin{align} \label{eq:channel_ope}
      A_{j-1}^\dagger \circ \mathcal{N}_{j-1}^{'}\circ A_{j-1}(\rho)&= \sum_{m = 0}^3 p_m (A_{j-1}^\dagger \sigma_{m} A_{j-1}) \rho (A_{j-1}^\dagger \sigma_{m} A_{j-1}).
   \end{align}
   Now we show that $L_j A_{j}^\dagger$ is one of the stabilizers of state $D_j^\dagger \ket{+}^{\otimes n}$ by identity below:
   \begin{align}
       L_j A_{j}^\dagger D_j^\dagger \ket{+}^{\otimes n} &=  L_j  (\text{Cl}_{j+1} \cdots\text{Cl}_d) L_d^\dagger (G_d \cdots G_{j+1}) D_j^\dagger \ket{+}^{\otimes n} \\
       &= L_j  (\text{Cl}_{j+1} \cdots\text{Cl}_d) L_d^\dagger D_d^\dagger \ket{+}^{\otimes n} \\
       &= L_j  (\text{Cl}_j^\dagger \cdots\text{Cl}_{1}^\dagger)  \ket{+}^{\otimes n} \\
       &= D_j^\dagger \ket{+}^{\otimes n}
   \end{align}
   Based on this, $(L_{j-1}A_{j-1}^\dagger)$ and $(A_{j-1}L_{j-1}^\dagger)$ are both Pauli operators. we decompose the Pauli channel as
   \begin{align}
       (A_{j-1}^\dagger \sigma_{m} A_{j-1}) &=L_{j-1}^\dagger (L_{j-1}A_{j-1}^\dagger)  \sigma_{m} (A_{j-1}L_{j-1}^\dagger) L_{j-1}\\
       &=(-1)^{f(m)}(L_{j-1}^\dagger \sigma_{m}L_{j-1}) ,
   \end{align}
The sign $(-1)^{f(m)}$ depends on the commutation property. Fortunately, this sign will fade out in channels operations in Eq.~\eqref{eq:channel_ope}, eventually gives the operation
\begin{align}
      A_{j-1}^\dagger\circ \mathcal{N}_{j-1}^{'}\circ A_{j-1}(\rho) &= \sum_{m = 0}^3 p_m (L_{j-1}^\dagger \sigma_{m}L_{j-1})  \rho (L_{j-1}^\dagger \sigma_{m}L_{j-1})\\
      &=L_{j-1}^\dagger\circ \mathcal{N}_{j-1}^{'}\circ L_{j-1}(\rho) \\
      &=L_{j-1}^\dagger\circ L_{j-1} \circ \mathcal{N}_{j-1}\circ L_{j-1}^\dagger \circ L_{j-1}(\rho) \\
      &= \mathcal{N}_{j-1}(\rho).
   \end{align}
   This gives exactly the original NISQ T-gadget circuit.
\end{proof}
In the end, the whole operator $\hat O$ is composed of phase gates $G_i, \prod CZ_i$, the shifted local Pauli noise $\mathcal{N}_i'$ and a local Clifford gate $L_d$ attached at the beginning of the input qubits, as shown in Fig.~\ref{fig:Percolated_stabilizer_state}(b) and equation below
\begin{align}
\label{eq:conjugated_amplitude_diagonal}
    \hat O=  (\mathcal{N'}_1^\dagger\circ \prod\mathcal{CZ}_1^\dagger\circ \mathcal{G}_1^\dagger)\circ \cdots\circ(\mathcal{N'}_d^\dagger\circ \prod\mathcal{CZ}_d^\dagger\circ \mathcal{G}_d^\dagger)\circ \mathcal{N'}_0^\dagger\left[ \ket{+}\!\bra{+}^{\otimes (n+t)} \right].
\end{align}
%%%%%%%%%%%%%%%%%%%%%%%%
%%%%%%%%%%%%%%%%%%%%%%%%%%%%%
%%%%%%%%%%%%%%%%%%%%%%%%%%%%%%%%%%

\subsubsection{Simulate incoherent noise through vertex deletion}
In this section, we show how to classical simulate the noise layer in this operator $\hat O$  through classical sampling graphs.
For unital incoherent noise, $\mathcal{N}$ is single qubit Pauli noise and $\mathcal{N}^\dagger$ is a valid qunatum channel.
This is equivalent to the evolution of IQP-circuits under local Pauli noise.

Many methods have been proposed to discuss the percolation behavior in the IQP-circuits~\cite{nelson2024polynomial,rajakumar2024polynomialtimeclassicalsimulationnoisy}. These methods are based on the components of $Z$-projective channel $\Lambda_z(\rho) = ({1}/{2})\rho+({1}/{2})Z\rho Z$ in Pauli noise that decoheres $\ket{+}$ state into uniform state ensemble $\{\ket{0}\!\bra{0}, \ket{1}\!\bra{1}\}$ , which will not entangled with the other $\ket{+}$  states under the following operation of phase gates. To be precise, we refer the theorem proved by Joel \emph{et.al.}~\cite{rajakumar2025polynomial}. 
\begin{lemma} \label{lemma:z-projective_of_Paulinoise}
    Any Pauli noise channel $\mathcal{N}_i'(\rho) = (1 - p_x-p_y-p_z)\rho +p_x X\rho X + p_y Y\rho Y + p_z Z\rho Z$ in Fig.~\ref{fig:Percolated_stabilizer_state}(b) can be decomposed as
    \begin{align}
        \mathcal{N}_{i}'(\rho) = (1 - 2p)\mathcal{N}_i^{(1)}(\rho) + 2p\mathcal{N}_i^{(2)} \circ \Lambda_z(\rho),
    \end{align}
    where $p = p_Z + \min(p_X,p_Y)$ , $N^{(2)}_i$ is Pauli noise channel with parameter $(0,\frac{\min(p_x, p_y)}{p},0)$  , and  $N^{(1)}_i$ is Pauli noise channel with parameter $(0,\frac{|p_x, p_y|}{1 - 2p},0)$ if $p_x < p_y$ and $(\frac{|p_x, p_y|}{1 - 2p},0,0)$ if $p_x \geq p_y$ .
\end{lemma}
Then each unital noise layer can be simulated through randomly remove this graph vertex with probability $2p$ and implementation of Pauli noise channel $\mathcal{N}_i^{(1)},\mathcal{N}_i^{(2)}$.

For non-unital nosie,  $\mathcal{N}^\dagger$ is not a trace preserving map and the operator $\hat O$ is not necessary a valid density matrix. 
For example, the conjugation of amplitude damping noise gives
\begin{align}
    \Lambda^\dagger(\frac{\mathbb{I} + a_1 \sigma_1 + a_2 \sigma_2 + a_3 \sigma_3}{2}) = \frac{(1 - a_3p)\mathbb{I} +\sqrt{1- p}\left( a_1 \sigma_1 + a_2 \sigma_2\right) + (1 - p)a_3 \sigma_3}{2},
\end{align}
which increases the trace from $1$ into $1 + p$. However, it is still possible to simulate single qubit noise through graph vertex deletion, with a price that we have to sample a "negative" density matrix $\{-\ket{0}\!\bra{0}, -\ket{1}\!\bra{1}\}$.
This is shown by the theorem below.

\begin{theorem}
\label{theo:sampling_nonunital_noise}
    Given a stabilizer state $S:=\ket{S}\!\bra{S}$ and an single qubit incoherent noise channel $\Lambda$ with parameters
 \begin{align}
    \mathcal{N}\left(\frac{1}{2}[\mathbb{I} + \mathbf{w} \cdot \vec \sigma]\right) = \frac{1}{2}[\mathbb{I} + (\mathbf{t} + \Gamma\mathbf{w}) \cdot  \vec\sigma] .
 \end{align}
Then the operator  $\mathcal{N}^\dagger(S)$ can be simulated by sampling signed and measured stabilizer states:
\begin{align}
    \{\pm S, \pm \sigma_i S \sigma_i^\dagger,\pm \ket{0}\!\bra{0} \otimes R_i^{\pm}, \pm\ket{1}\!\bra{1} \otimes R_i^{\pm}\}
\end{align}
with different probability. Here $R_i^{\pm}$ is stabilizer state through measuring the qubit in stabilzier state $\ket{S}$ on $(\mathbb{I} \pm \sigma_i)/2$ basis.
\end{theorem}
\begin{proof}
    Without loss of generality, we can suppose the single qubit noise channel is implemented on the first qubit and the state satisfied
    \begin{align}
        S:=\ket{S}\!\bra{S} = \mathbb{I} \otimes Q_0 + \sum_i \sigma_i\otimes Q_i,
    \end{align}
    where $Q_i$ is linear combination of Pauli string that support on the $2,\cdots,n$ qubits. 
By definition, the measurement of $\sigma_i \otimes\mathbb{I} \otimes \cdots \otimes \mathbb{I}$ and post selection on result $\{+,-\}$ will give the operator $M_i^{\pm}$. Thus we conclude
\begin{align}
    M_i^{\pm} &=[(\mathbb{I} \pm \sigma_i)/2]S[(\mathbb{I} \pm \sigma_i)/2]  \\
    &= [(\mathbb{I} \pm \sigma_i)/2]\left(\mathbb{I} \otimes Q_0 + \sum_j \sigma_j\otimes Q_j \right)[(\mathbb{I} \pm \sigma_i)/2]\\
     &= [(\mathbb{I} \pm \sigma_i)/2]\left(\mathbb{I} \otimes Q_0 +  \sigma_i\otimes Q_i \right)\\
     &=\frac{\mathbb{I} \pm \sigma_i}{2}\left(\mathbb{I} \otimes Q_0 \pm  \sigma_i\otimes Q_i \right)
\end{align}
Noted that if $\ket{S}$ is a stabilizer state and the first qubit is entangled with the rest of the system, $M_i^{\pm}$ also has the form
\begin{align}
    M_i^{\pm} &= (\frac{\mathbb{I} \pm \sigma_i}{2}) \otimes \frac{1}{2}\ket{R^{\pm}_i}\!\bra{R^{\pm}_i},
\end{align}
where $\ket{R_i^\pm}$ is another normalized stabilizer state on $\{2,\cdots,n\}$ qubits and has lower entanglement width compare to original stabilizer state $\ket{S}$.
Then we conclude that
\begin{align}
   &\ket{R^{+}_i}\!\bra{R^{+}_i} =2Q_0+ 2Q_i ,\\
    & \ket{R^{-}_i}\!\bra{R^{-}_i} =2Q_0- 2Q_i .
\end{align}
For unital incoherent noise, supposed $\Gamma_1 > \Gamma_2$, otherwise we exchanged the charator of $1,2$ in the following proof. It can be conclude that
\begin{align}
    \Lambda^\dagger\left(\ket{S}\!\bra{S}\right) &= \mathbb{I} \otimes Q_{0} + \sum_{i } \Gamma_i \sigma_i \otimes Q_i \\
    &=\frac{\Gamma_1+ \Gamma_2}{2} \left( \mathbb{I} \otimes Q_0 + \sigma_1\otimes Q_1 + \sigma_2\otimes Q_2 + \sigma_3\otimes Q_3\right)\\
     &\quad\quad \quad +\frac{\Gamma_1- \Gamma_2}{2} \left( \mathbb{I} \otimes Q_0 + \sigma_1\otimes Q_1 - \sigma_2\otimes Q_2 - \sigma_3\otimes Q_3\right)\\
   & \quad \quad \quad \quad  +(\Gamma_3 - \Gamma_2)\sigma_3 \otimes Q_3 + (1 - \Gamma_1)\mathbb{I} \otimes Q_0\\
    &= \frac{\Gamma_1+ \Gamma_2}{2}  S + \frac{\Gamma_1- \Gamma_2}{2} \sigma_1 S \sigma_1 + (\Gamma_3 - \Gamma_2)\sigma_3 \otimes(\frac{R_3^+ - R_3^-}{4})  + (1- \Gamma_1)\mathbb{I} \otimes (\frac{R_3^+ + R_3^-}{4}) \\
    &= \frac{\Gamma_1+ \Gamma_2}{2}  S + \frac{\Gamma_1- \Gamma_2}{2} \sigma_1 S \sigma_1 \\
    &\quad \quad + \frac{ 1 - \Gamma_1- \Gamma_2+\Gamma_3 }{4}\frac{\mathbb{I} + \sigma_3 }{2} \otimes {R_3^+}+\frac{1 - \Gamma_1 + \Gamma_2-\Gamma_3}{4}\frac{\mathbb{I}-\sigma_3}{2} \otimes {R_3^+}\\
    &\quad \quad + \frac{ 1 - \Gamma_1- \Gamma_2+\Gamma_3 }{4}\frac{\mathbb{I} - \sigma_3 }{2} \otimes {R_3^-}+\frac{1 - \Gamma_1 + \Gamma_2-\Gamma_3}{4}\frac{\mathbb{I}+\sigma_3}{2} \otimes {R_3^-}.
\end{align}
For a valid Pauli channel, we conclude that $\Gamma_1 = 1 - 2p_2 - 2p_3,\Gamma_2 = 1 - 2p_1 - 2p_3$ and $\Gamma_3 = 1 - 2p_1 - 2p_2$. Then $1 + \Gamma_3 = 2 - 2p_1 - 2p_2 >2 - 2p_1 - 2p_2 - 4p_3 = \Gamma_1 + \Gamma_2$ should always be satisfied. The back-propagation can be simulated through samping stabilizer state $\{\ket{S}, \sigma_1 \ket{S}, \ket{\pm} \otimes \ket{R_3^{\pm}}\}$ with corresponding propability. The propability that the entanglement structure is unchanged is $\Gamma_1$. This is coincident with previous result~\cite{rajakumar2025polynomial}.

For non-unital noise, the back-propagation result writes 
\begin{align}
    \Lambda^\dagger\left(\ket{S}\!\bra{S}\right) = \mathbb{I} \otimes Q_{0} + \sum_{i } \Gamma_i \sigma_i \otimes Q_i + \sum_i t_i \mathbb{I} \otimes Q_i,
\end{align} 
with an additional term $ \mathbb{I} \otimes \sum_i t_iQ_i$ added.
Again we supposed  $\Gamma_1 > \Gamma_2$ and notice that $\sum_i t_i Q_i = - (\sum_i t_i)Q_0 + (1/2) \sum_i t_i R_i^+$.
, the back-propagation result can be decomposed similarly as
\begin{align}
    \Lambda^\dagger\left(\ket{S}\!\bra{S}\right) &= \mathbb{I} \otimes Q_{0} + \sum_{i } \Gamma_i \sigma_i \otimes Q_i + \sum_i t_i \mathbb{I} \otimes Q_i\\
    &= \frac{\Gamma_1+ \Gamma_2}{2}  S + \frac{\Gamma_1- \Gamma_2}{2} \sigma_1 S \sigma_1 +(\Gamma_3 - \Gamma_2)\sigma_3 \otimes Q_3 \\
    &\quad \quad + (1 - \Gamma_1-\sum_i t_i)\mathbb{I} \otimes Q_0 + \frac{1}{2}\mathbb{I} \otimes \sum_i t_i R_i^+\\
     &= \frac{\Gamma_1+ \Gamma_2}{2}  S + \frac{\Gamma_1- \Gamma_2}{2} \sigma_1 S \sigma_1 \\
    & \quad + \frac{ 1-\sum_i t_i - \Gamma_1- \Gamma_2+\Gamma_3 }{4}\frac{\mathbb{I} + \sigma_3 }{2} \otimes {R_3^+}+\frac{1 -\sum_i t_i- \Gamma_1 + \Gamma_2-\Gamma_3}{4}\frac{\mathbb{I}-\sigma_3}{2} \otimes {R_3^+}\\
    &\quad + \frac{ 1 -\sum_i t_i- \Gamma_1- \Gamma_2+\Gamma_3 }{4}\frac{\mathbb{I} - \sigma_3 }{2} \otimes {R_3^-}+\frac{1 -\sum_i t_i- \Gamma_1 + \Gamma_2-\Gamma_3}{4}\frac{\mathbb{I}+\sigma_3}{2} \otimes {R_3^-}\\
    & \quad +\sum_i\frac{t_i}{2}\frac{\mathbb{I} - \sigma_3}{2}\otimes  R_i^+ +\sum_i\frac{t_i}{2}\frac{\mathbb{I} + \sigma_3}{2}\otimes  R_i^+.
\end{align}
We should change the signal of single qubit density matrix ${(\mathbb{I} \pm \sigma_3) }/{2}$ if ${( 1-\sum_i t_i - \Gamma_1- \Gamma_2+\Gamma_3 )}/{4} <0$ or ${( 1-\sum_i t_i - \Gamma_1+ \Gamma_2-\Gamma_3 )}/{4} <0$.
Specially, for amplitude damping noise, $\Gamma_1 = \Gamma_2 = \sqrt{1 - p}, \Gamma_3 = 1-p$ and $t_1 = t_2 = 0, t_3 = -p$. Then we conclude that:
\begin{align}
    \Lambda^\dagger\left(\ket{S}\!\bra{S}\right) 
     &=  \sqrt{1 - p}  S  \\
    & \quad + \frac{\sqrt{1 - p} -  (1-p) }{2}\left(-\frac{\mathbb{I} + \sigma_3 }{2} \right)\otimes {R_3^+}\\
    &\quad + \frac{ \sqrt{1 - p} -(1-p) }{2}\left(-\frac{\mathbb{I} - \sigma_3 }{2}\right) \otimes {R_3^-}\\
    & \quad +\frac{p}{2}\frac{\mathbb{I} - \sigma_3}{2}\otimes  R_3^+ +\frac{p}{2}\frac{\mathbb{I} + \sigma_3}{2}\otimes  R_3^+.
\end{align}
Noted that the removed qubit are distributed with a "negative" quantum state $-\left({\mathbb{I} \pm \sigma_3 }/{2}\right)$. This is not a valid quantum state, but they will keep disentangled with the other state through the $CZ$ operation. Thus it is possible to trace these single qubit "negative" state.  
\end{proof}


The mixed stabilizer state that overlaps with product state $\ket{T}^{\otimes t}$ is $\sigma = \text{tr}_{[n]} \left( \rho \ket{0}\!\bra{0}^{\otimes n}\right)$. 
$\sigma$ has high probability generates percolated stabilizer states, thus its overlap between $\ket{T}^{\otimes t}$ can be calculated Based on Lemma~\ref{lemma: stabilizer_to_graph}, all the Clifford gates in the circuit can be reduced to equivalent phase gates, as explained in the theorem~\ref{lemma:Noisy_reduction}.

Further more, using theorem~\ref{lemma:Noisy_reduction} , all the Clifford gates in T-gadget circuit can be equivalently transformed into phase gates and by-product local Clifford gates $L$ on input qubits, as illustrated in Fig.~\ref{fig:Percolated_stabilizer_state}(d). Those local Clifford gates can move back to the beginning of the circuits and shift Pauli noise $\mathcal{N}_i$ to $\mathcal{N}_{i}'$, which is also a Pauli noise whose distribution can be efficiently computed. Local Clifford gates also transform T-gadget $CZ$ gates into control-Pauli gates with target on input qubits. 



With the help of lemma~\ref{lemma:z-projective_of_Paulinoise}, the back-propagation of  $\ket{+}^{\otimes (n+t)}$ state can be fully described by operation of classical graph. The gates $G_i$ change the edge of the graph. The Pauli local noise randomly removes input-qubit vertexes while the T-gadgets $CZ$ gates introduce ancilla-qubit vertex and connect them with multiple input-qubit vertexes. When the graph propagates to the beginning of the circuit, the local-Clifford gates $L$ and circuit input state $\ket{0}^{\otimes n}$ introduce Pauli projections on this graph states. This also equivalent to certain removing rule for vertexes of the graph, as illustrate by the lemma below~\cite{mor2023noisy}.
\begin{lemma}\label{lemma:Pauli_measurement_on_graph}
    The graph state $\ket{G}$ under goes a Pauli measurement at vertex $a$, resulting in a new graph state $\ket{G'}$ that is locally Clifford equivalent. Specifically:
\begin{enumerate}
    \item After a $Z$-measurement, the graph is modified by deleting all edges incident to vertex $a$.
    \item For a $Y$-measurement, the neighborhood of $a$ is flipped, and then all edges incident to $a$ are removed.
    \item For an $X$-measurement, choose a neighboring vertex $b$ of $a$. First, flip the neighborhood of $b$, then flip the neighborhood of $a$, followed by removing the edges incident to $a$, and finally flip the neighborhood of $b$ again.
\end{enumerate}
\end{lemma}
Going through the rule in lemma~\ref{lemma:Pauli_measurement_on_graph} finally gives random graph states along with locally Clifford gates induced by the Pauli measurements.
If the graph of the graph state is percolated into fraction smaller than $\mathcal{O}(\ln n)$ vertex, we can rigorously calculate its inner product with the magic states $(H\ket{T})^{\otimes n}$  fraction-by-fraction.   
We prove that the value of inner products is bound by $\mathcal{O}(2^{-t/2})$, thus bounds the variance of Monte-Carlo sampling method.
\begin{theorem}\label{theo:variance_of_percolation_sampling}
    The algorithm samples from graph states along with local Clifford gates $\{\ket{G_i},L_i\}_i$ ,then each of them satisfied $|\bra{G_i}L_i(H\ket{T})^{\otimes n}|^2 < 2^{-t}$.   
\end{theorem}
\begin{proof}
    Consider one of the qubits with $H\ket{T}$ inputs, its post-selection writes 
    \begin{align}
       \bra{+}_t CP_{t,n}\left[(S^a Z^b H\ket{T}_t) \otimes \ket{\phi}_n\right] &=\bra{+}_t CQ_{t,n}\left[( H\ket{T}_t) \otimes \ket{\phi}_n\right]\\
       &=\bra{+}_t\left[\cos(\frac{\pi}{8}) \ket{0}_t \ket{\phi}_n+i\sin(\frac{\pi}{8}) \ket{1}_t (Q\ket{\phi})_n\right]\\
       &=\frac{\ket{+}_t}{\sqrt{2}} e^{-i\frac{\pi}{8}Q} \ket{\phi}
    \end{align}
    $Q$ is another Pauli operator. This is because $S$ gate on target qubit can always equivalent to control-$Y$ 
 gates while $Z$ gates is equivalent to control-$(-I)$ gates. This means that even if we remove some $CZ$ gates due to the effect of noise, the circuit left are still a legal gadget circuits with $t$ gadget implementation, which still has the post-selection overhead $2^{-t/2}$. 
\end{proof}
The Hoeffding inequalities show that for random variables bounded in the interval $(-2^{-t},2^{-t})$, $N$ samples will estimate its mean within additive error $2^{-t}\varepsilon$ of the actual mean with probability exceeding $1-\delta$ where $\delta = 2 \cdot \text{exp}{(-2N(2^{-t}\varepsilon)^2/(4 \times2^{-2t}))}$. In other  words, the desired accuracy is guaranteed by using 
\begin{align}
    N = \frac{2}{\varepsilon^2} \ln \frac{2}{\delta}
\end{align}
samples. Because the result of the original circuit is $2^t$ times the amplitude of T-gadget circuit, estimating T-gadget amplitude within additive error $2^{-t}\varepsilon$ gives the amplitude of the original circuit with $\varepsilon$ additive error. 

The purity of $\sigma $ is bounded by that of $\rho$, which concludes
\begin{align}
    \text{tr}(\sigma^2) =2^{t} \langle \!\langle \sigma|\sigma\rangle\!\rangle\leq 2^{t} \langle \!\langle \rho|\rho\rangle\!\rangle = 2^{-n}\sum_{E_0,E_i} 2^{-n-t}(1 - \gamma)^{2f(E_0,E_i)}
\end{align}
If the term $\sum_{E_0,E_i} 2^{-n-t}(1 - \gamma)^{2f(E_0,E_i)}$ has $\text{poly}(n)$ scale, we concludes that $\text{tr}(\sigma^2) \leq {\text{Poly}(n)}{2^{-n}}$. 

\begin{theorem}[Percolation on low purity stabilizer states]
    Given a mixed state of stabilizers. If $\text{tr}(\sigma^2) < {\text{Poly}(n)}{2^{-n}}$, There exist one convex decomposition into pure states $\sigma= \sum_i p_i \ket{\psi_i}\!\bra{\psi_i}, \ket{\psi_i} \in \text{Stab}_n$ such that each $ \ket{\psi_i}$ will localized on $\mathcal{O}(\ln n)$ number of qubits.    
\end{theorem}

The simple structure of graph states also gives more powerful testing strategy compared to those in the  IQP-circuits~\cite{nelson2024polynomial,rajakumar2024polynomialtimeclassicalsimulationnoisy}. It has been shown that if the tree width of graph $G$ is $O(\ln n)$, with $n$ the vertex numbers, it turns out that there exist a polynomial algorithm that calculates the inner product between $n$-qubit graph state $\ket{G}$ and arbitrary product states $\ket{\nu}^{\otimes n}$~\cite{van2007classical}.


\subsection{Graphical compilation for weak-simulation}
As illustrated in the last section, we can calculate the sampling amplitude through calculating the inner product between magic inputs and stabilizer state. It is natual to ask, whether we can fully replace the quantum circuit with a measurement on graph state. In this section, we show that this is possible and gives a concrete algortihm that compiles a given Clifford+T circuit into POVM on certain graph state.
 
We first consider a noiseless circuit composed of universal gate set $\{CZ, T,H\}$ . Note that we can teleport all the $T$ gates, move $H$ gates to the end of the circuit and replace control-$X$ gates with multiple control-$Z$ gates. A concrete example for this process is illustrated in Figure below.

\begin{figure}[!htbp]
    \centering
    \begin{quantikz}[row sep=0.2cm, column sep=0.1cm]
  \lstick{$\ket{+}$}   \qw &\ctrl{1}       &\qw         &\gate{T}  &\gate{H}      &\ctrl{1}    & \qw       &\qw     &\gate{Z^{m_1}}&\meter{\ket{+}}\\
  \lstick{$\ket{+}$}   \qw & \control{}    &\ctrl{1}    &\gate{T}  &\qw     & \control{} & \ctrl{1}  &\qw          &\gate{Z^{m_2}}&\meter{}\\
  \lstick{$\ket{+}$}   \qw &\qw            & \control{}      &\gate{T}  &\gate{H} &\qw         &\control{} &\gate{T}&\gate{Z^{m_3}}&\meter{}
\end{quantikz}
       \begin{quantikz}[row sep=0.2cm, column sep=0.1cm]
  \lstick{$\ket{+}$}   \qw &\ctrl{1}      &\qw            &\qw      &\ctrl{3}  &\qw        &\qw        &\gate{H}      &\ctrl{1}    & \qw       &\qw        &\qw      &\meter{\ket{+}}\\
  \lstick{$\ket{+}$}   \qw &\control{}    &\ctrl{1}       &\qw      &\qw       &\ctrl{3}   &\qw        &\qw & \control{} & \ctrl{1}  &\qw        &\qw      &\meter{}\\
  \lstick{$\ket{+}$}   \qw &\qw           & \control{}    &\qw      &\qw       &\qw        &\ctrl{3}   &\gate{H} &\qw         &\control{} &\ctrl{4}   &\qw     &\meter{}\\
  \lstick{$\ket{A}$}   \qw &\qw           &\qw            &\qw      &\control{}&\qw        &\qw        & \qw     &\qw         &\qw        &\qw        &\qw      &\meter{}\\
  \lstick{$\ket{A}$}   \qw &\qw           &\qw            &\qw      &\qw       &\control{} &\qw        & \qw     &\qw         &\qw        &\qw        &\qw      &\meter{}\\
  \lstick{$\ket{A}$}   \qw &\qw           &\qw            &\qw      &\qw       &\qw        &\control{} &\qw      &\qw         &\qw        &\qw        &\qw      &\meter{}       \\
  \lstick{$\ket{A}$}   \qw &\qw           &\qw            &\qw      &\qw       &\qw        &\qw        &\qw      &\qw         &\qw        &\control{} &\qw      &\meter{}
\end{quantikz}
       \begin{quantikz}[row sep=0.2cm, column sep=0.1cm]
  \lstick{$\ket{+}$}   \qw   &\gate{H} &\targ{}     &\qw             &\targ{}  &\qw        &\qw         &\ctrl{1}   & \qw       &\qw          &\qw      &\meter{\ket{+}}\\
  \lstick{$\ket{+}$}   \qw &\qw        &\ctrl{-1}     &\ctrl{1}     &\qw       &\control{} &\qw         & \control{}& \ctrl{1}  &\qw          &\qw      &\meter{}\\
  \lstick{$\ket{+}$}   \qw  &\gate{H}  &\qw           & \targ{}    &\qw       &\qw        &\targ{}    &\qw         &\control{} &\ctrl{4}    &\qw      &\meter{}\\
  \lstick{$\ket{A}$}   \qw  & \qw      &\qw           &\qw         &\ctrl{-3}  &\qw        &\qw         &\qw        &\qw        &\qw          &\qw      &\meter{}\\
  \lstick{$\ket{A}$}   \qw  & \qw      &\qw           &\qw         &\qw       &\ctrl{-3}   &\qw         &\qw       &\qw        &\qw           &\qw      &\meter{}\\
  \lstick{$\ket{A}$}   \qw  &\qw       &\qw           &\qw         &\qw       &\qw        &\ctrl{-3}  &\qw         &\qw        &\qw         &\qw      &\meter{}       \\
  \lstick{$\ket{A}$}   \qw  &\qw       &\qw           &\qw         &\qw       &\qw        &\qw         &\qw        &\qw        &\control{}   &\qw      &\meter{}
\end{quantikz}
\\
%        \begin{quantikz}[row sep=0.2cm, column sep=0.1cm]
%   \lstick{$\ket{+}$}   \qw &\gate{H} &\targ{}         &\qw             &\targ{}   &\qw       &\qw       &\qw        &\qw         &\ctrl{1}   & \qw       &\qw         &\qw      &\meter{\ket{+}}\\
%   \lstick{$\ket{+}$}   \qw &\qw      &\ctrl{-1}       &\ctrl{1}     &\qw        &\control{}&\qw       &\qw        &\control{}  &\control{} & \ctrl{1}  &\qw         &\qw      &\meter{}\\
%   \lstick{$\ket{+}$}   \qw &\gate{H} &\qw             & \targ{}    &\qw        &\qw        &\qw       &\qw        &\qw         &\qw        &\control{} &\ctrl{4}    &\qw      &\meter{}\\
%   \lstick{$\ket{A}$}   \qw & \qw     &\qw             &\qw         &\ctrl{-3}   &\qw       &\qw       &\qw        &\qw         &\qw        &\qw        &\qw         &\qw      &\meter{}\\
%   \lstick{$\ket{A}$}   \qw & \qw     &\qw             &\qw         &\qw        &\ctrl{-3}  &\ctrl{1}  &\qw        &\qw         &\qw        &\qw        &\qw         &\qw      &\meter{}\\
%   \lstick{$\ket{A}$}   \qw &\qw      &\qw             &\qw         &\qw        &\qw        &\control{}&\ctrl{1}   &\ctrl{-4}   &\qw        &\qw        &\qw         &\qw      &\meter{}       \\
%   \lstick{$\ket{A}$}   \qw &\qw      &\qw             &\qw         &\qw        &\qw        &\qw       &\control{} &\qw         &\qw        &\qw        &\control{}  &\qw      &\meter{}
% \end{quantikz}
%        \begin{quantikz}[row sep=0.2cm, column sep=0.1cm]
%   \lstick{$\ket{+}$}   \qw &\gate{H} &\targ{}         &\qw             &\qw       &\qw       &\qw       &\qw        &\qw         &\ctrl{1}   & \qw       &\qw         &\qw      &\meter{\ket{+}}\\
%   \lstick{$\ket{+}$}   \qw &\qw      &\ctrl{-1}       &\ctrl{1}     &\control{}  &\control{}&\qw       &\qw        &\control{}  &\control{} & \ctrl{1}  &\qw         &\qw      &\meter{}\\
%   \lstick{$\ket{+}$}   \qw &\gate{H} &\qw             & \targ{}    &\qw        &\qw        &\qw       &\qw        &\qw         &\qw        &\control{} &\ctrl{4}    &\qw      &\meter{}\\
%   \lstick{$\ket{A}$}   \qw & \qw     &\qw             &\qw         &\ctrl{-2}   &\qw       &\qw       &\qw        &\qw         &\qw        &\qw        &\qw         &\qw      &\meter{}\\
%   \lstick{$\ket{A}$}   \qw & \qw     &\qw             &\qw         &\qw        &\ctrl{-3}  &\ctrl{1}  &\qw        &\qw         &\qw        &\qw        &\qw         &\qw      &\meter{}\\
%   \lstick{$\ket{A}$}   \qw &\qw      &\qw             &\qw         &\qw        &\qw        &\control{}&\ctrl{1}   &\ctrl{-4}   &\qw        &\qw        &\qw         &\qw      &\meter{}       \\
%   \lstick{$\ket{A}$}   \qw &\qw      &\qw             &\qw         &\qw        &\qw        &\qw       &\control{} &\qw         &\qw        &\qw        &\control{}  &\qw      &\meter{}
% \end{quantikz}
%        \begin{quantikz}[row sep=0.2cm, column sep=0.1cm]
%   \lstick{$\ket{+}$}   \qw &\gate{H}   &\qw            &\qw       &\qw       &\qw       &\qw        &\qw         &\ctrl{1}   & \qw       &\qw         &\qw      &\meter{\ket{+}}\\
%   \lstick{$\ket{+}$}   \qw &\qw        &\ctrl{5}      &\control{}  &\control{}&\qw       &\qw        &\control{}  &\control{} & \ctrl{1}  &\qw         &\qw      &\meter{}\\
%   \lstick{$\ket{+}$}   \qw &\gate{H}   &\qw          &\qw        &\qw        &\qw       &\qw        &\qw         &\qw        &\control{} &\ctrl{4}    &\qw      &\meter{}\\
%   \lstick{$\ket{A}$}   \qw & \qw       &\qw          &\ctrl{-2}   &\qw       &\qw       &\qw        &\qw         &\qw        &\qw        &\qw         &\qw      &\meter{}\\
%   \lstick{$\ket{A}$}   \qw & \qw       &\qw          &\qw        &\ctrl{-3}  &\ctrl{1}  &\qw        &\qw         &\qw        &\qw        &\qw         &\qw      &\meter{}\\
%   \lstick{$\ket{A}$}   \qw &\qw        &\qw          &\qw        &\qw        &\control{}&\ctrl{1}   &\ctrl{-4}   &\qw        &\qw        &\qw         &\qw      &\meter{}       \\
%   \lstick{$\ket{A}$}   \qw &\qw        &\control{}   &\qw        &\qw        &\qw       &\control{} &\qw         &\qw        &\qw        &\control{}  &\qw      &\meter{}
% \end{quantikz}

    \caption{ In the embedding setting, the equivalent qubit gates implement the hyper control-\(X\) gate with \(d = 4\).}
    \label{fig:comp_1}
\end{figure}


\subsection{Hypergraph and weighted graph compilation}

\end{document}